% problem_id: S001-lemma-higher-direct-image
% source_id: S001 (Stacks Project)
% source_file: spaces-cohomology.tex
% statement_locator: spaces-cohomology.tex lines 121--127
% proof_locator: spaces-cohomology.tex lines 129--235
% env: lemma
% pairing: tight (verified)
% status: complete (statement + author proof verbatim + reason + locators)
%
\begin{lemma}
\label{lemma-higher-direct-image}
Let $S$ be a scheme. Let $f : X \to Y$ be a morphism of algebraic spaces
over $S$. If $f$ is quasi-compact and quasi-separated, then $R^if_*$
transforms quasi-coherent $\mathcal{O}_X$-modules into
quasi-coherent $\mathcal{O}_Y$-modules.
\end{lemma}

\begin{proof}
Let $V \to Y$ be an \'etale morphism where $V$ is an affine scheme. Set
$U = V \times_Y X$ and denote $f' : U \to V$ the induced morphism.
Let $\mathcal{F}$ be a quasi-coherent $\mathcal{O}_X$-module. By
Properties of Spaces, Lemma
\ref{spaces-properties-lemma-pushforward-etale-base-change-modules}
we have
$R^if'_*(\mathcal{F}|_U) = (R^if_*\mathcal{F})|_V$.
Since the property of being a quasi-coherent module is local in the
\'etale topology on $Y$ (see
Properties of Spaces, Lemma
\ref{spaces-properties-lemma-characterize-quasi-coherent})
we may replace $Y$ by $V$, i.e., we may assume $Y$ is an affine scheme.

\medskip\noindent
Assume $Y$ is affine. Since $f$ is quasi-compact we see that $X$
is quasi-compact. Thus we may choose an affine scheme $U$ and a surjective
\'etale morphism $g : U \to X$, see
Properties of Spaces,
Lemma \ref{spaces-properties-lemma-quasi-compact-affine-cover}.
Picture
$$
\xymatrix{
U \ar[r]_g \ar[rd]_{f \circ g} & X \ar[d]^f \\
& Y
}
$$
The morphism $g : U \to X$ is representable, separated
and quasi-compact because $X$ is quasi-separated. Hence the lemma
holds for $g$ (by the discussion above the lemma).
It also holds for $f \circ g : U \to Y$ (as this is a morphism
of affine schemes).

\medskip\noindent
In the situation described in the previous paragraph we will show by
induction on $n$ that $IH_n$: for any quasi-coherent sheaf $\mathcal{F}$
on $X$ the sheaves $R^if\mathcal{F}$
are quasi-coherent for $i \leq n$.
The case $n = 0$ follows from
Morphisms of Spaces, Lemma \ref{spaces-morphisms-lemma-pushforward}.
Assume $IH_n$. In the rest of the proof we show that $IH_{n + 1}$ holds.

\medskip\noindent
Let $\mathcal{H}$ be a quasi-coherent $\mathcal{O}_U$-module.
Consider the Leray spectral sequence
$$
E_2^{p, q} = R^pf_* R^qg_* \mathcal{H}
\Rightarrow
R^{p + q}(f \circ g)_*\mathcal{H}
$$
Cohomology on Sites, Lemma \ref{sites-cohomology-lemma-relative-Leray}.
As $R^qg_*\mathcal{H}$ is quasi-coherent by $IH_n$ all the sheaves
$R^pf_*R^qg_*\mathcal{H}$ are quasi-coherent for $p \leq n$.
The sheaves $R^{p + q}(f \circ g)_*\mathcal{H}$ are all
quasi-coherent (in fact zero for $p + q > 0$ but we do not need this).
Looking in degrees $\leq n + 1$ the only module which we do not
yet know is quasi-coherent is $E_2^{n + 1, 0} = R^{n + 1}f_*g_*\mathcal{H}$.
Moreover, the differentials
$d_r^{n + 1, 0} : E_r^{n + 1, 0} \to E_r^{n + 1 + r, 1 - r}$
are zero as the target is zero. Using that $\QCoh(\mathcal{O}_X)$
is a weak Serre subcategory of $\textit{Mod}(\mathcal{O}_X)$
(Properties of Spaces, Lemma
\ref{spaces-properties-lemma-properties-quasi-coherent}) it
follows that $R^{n + 1}f_*g_*\mathcal{H}$
is quasi-coherent (details omitted).

\medskip\noindent
Let $\mathcal{F}$ be a quasi-coherent $\mathcal{O}_X$-module.
Set $\mathcal{H} = g^*\mathcal{F}$. The adjunction mapping
$\mathcal{F} \to g_*g^*\mathcal{F} = g_*\mathcal{H}$ is injective
as $U \to X$ is surjective \'etale. Consider the exact sequence
$$
0 \to \mathcal{F} \to g_*\mathcal{H} \to \mathcal{G} \to 0
$$
where $\mathcal{G}$ is the cokernel of the first map and in particular
quasi-coherent. Applying the long exact cohomology sequence we obtain
$$
R^nf_*g_*\mathcal{H} \to
R^nf_*\mathcal{G} \to
R^{n + 1}f_*\mathcal{F} \to
R^{n + 1}f_*g_*\mathcal{H} \to
R^{n + 1}f_*\mathcal{G}
$$
The cokernel of the first arrow is quasi-coherent and
we have seen above that $R^{n + 1}f_*g_*\mathcal{H}$ is quasi-coherent.
Thus $R^{n + 1}f_*\mathcal{F}$ has a $2$-step filtration
where the first step is quasi-coherent and the second a submodule of
a quasi-coherent sheaf. Since $\mathcal{F}$ is an arbitrary quasi-coherent
$\mathcal{O}_X$-module, this result also holds for $\mathcal{G}$.
Thus we can choose an exact sequence
$0 \to \mathcal{A} \to R^{n + 1}f_*\mathcal{G} \to \mathcal{B}$
with $\mathcal{A}$, $\mathcal{B}$ quasi-coherent $\mathcal{O}_Y$-modules.
Then the kernel $\mathcal{K}$ of
$R^{n + 1}f_*g_*\mathcal{H} \to R^{n + 1}f_*\mathcal{G}
\to \mathcal{B}$ is quasi-coherent, whereupon we obtain a map
$\mathcal{K} \to \mathcal{A}$ whose kernel $\mathcal{K}'$ is
quasi-coherent too. Hence $R^{n + 1}f_*\mathcal{F}$ sits in an exact
sequence
$$
R^nf_*g_*\mathcal{H} \to
R^nf_*\mathcal{G} \to
R^{n + 1}f_*\mathcal{F} \to \mathcal{K}' \to 0
$$
with all modules quasi-coherent except for possibly $R^{n + 1}f_*\mathcal{F}$.
We conclude that $R^{n + 1}f_*\mathcal{F}$ is quasi-coherent, i.e.,
$IH_{n + 1}$ holds as desired.
\end{proof}
